\section{晚明的朝廷、法律与文化连结}

本节以朝廷程序、法律语言、家庭与性别、城市文化和海外联系为问题线索。正史、实录、律例与奏议通常更容易保存官员、男性精英和国家分类；女性、家庭成员、工匠、城市居民与海外中介则常仅在墓志、契约、地方志、碑刻、航海记录或案件中留下片段。材料本身的偏向，应当成为叙事的一部分。

\subsection{人事、审案与政治语言}

\eventanchor{ML-1628-004}{明·1628（崇祯元年）}
\paragraph{1628（崇祯元年）：朝鲜、辽东和海上关系的本年材料提供外部观察位置，明、后金（清）与地方势力的称谓需具体化。（材料核查重点：诏令或奏议的形成与发受链）} 此事把权力、法律或文化关系推进到可辨认的节点。账本的证据限度是来源导向的年度桥接条目：本年材料可定位相关政务线索；具体月日、发文机关、地域和执行结果仍须逐项回查，不能以制度文本或奏报替代实际效果。可资核查的材料为《崇祯长编》崇祯元年条及《明史·思宗本纪》；《明史·外国传》；《朝鲜王朝实录》、琉球《历代宝案》或海外档案。它们能够说明制度如何表述自身、某种命令或交涉如何被记录，却未必能直接复原家庭内部关系、性别分工、城市日常或海外行动者的意图。事件发生的条件涉及朝贡名义、边境安全与港口贸易的利益使交涉持续发生；其后留下的制度与社会后果为它为后续册封、互市或海防安排留下文书与跨政权比较线索。桥接条目只标示可回查的年度问题与材料链；实录有中央选择性，崇祯期尤其需要奏疏、地方、朝鲜及满文材料交叉。；本条特别核对诏令或奏议的形成与发受链。

\eventanchor{ML-1628-005}{明·1628（崇祯元年）}
\paragraph{1628（崇祯元年）：辽饷、剿饷、练饷及地方征解的本年文书显示财政负担，定额、征收与实际流向不得混为一谈。（材料核查重点：诏令或奏议的形成与发受链）} 此事把权力、法律或文化关系推进到可辨认的节点。账本的证据限度是来源导向的年度桥接条目：本年材料可定位相关政务线索；具体月日、发文机关、地域和执行结果仍须逐项回查，不能以制度文本或奏报替代实际效果。可资核查的材料为《崇祯长编》崇祯元年条及《明史·思宗本纪》；《明会典》相关条目；地方志、黄册/契约/河工材料。它们能够说明制度如何表述自身、某种命令或交涉如何被记录，却未必能直接复原家庭内部关系、性别分工、城市日常或海外行动者的意图。事件发生的条件涉及财政征解、交通工程或货币支付的压力使相关措施进入政务；其后留下的制度与社会后果为其法定安排不等于地方执行，后续须以地域材料追踪负担与成效。桥接条目只标示可回查的年度问题与材料链；实录有中央选择性，崇祯期尤其需要奏疏、地方、朝鲜及满文材料交叉。；本条特别核对诏令或奏议的形成与发受链。

\eventanchor{ML-1628-006}{明·1628（崇祯元年）}
\paragraph{1628（崇祯元年）：辽饷、剿饷、练饷及地方征解的本年文书显示财政负担，定额、征收与实际流向不得混为一谈。（材料核查重点：报称信息的来源、抄传与行政分类）} 此事把权力、法律或文化关系推进到可辨认的节点。账本的证据限度是来源导向的年度桥接条目：本年材料可定位相关政务线索；具体月日、发文机关、地域和执行结果仍须逐项回查，不能以制度文本或奏报替代实际效果。可资核查的材料为《崇祯长编》崇祯元年条及《明史·思宗本纪》；《明会典》相关条目；地方志、黄册/契约/河工材料。它们能够说明制度如何表述自身、某种命令或交涉如何被记录，却未必能直接复原家庭内部关系、性别分工、城市日常或海外行动者的意图。事件发生的条件涉及财政征解、交通工程或货币支付的压力使相关措施进入政务；其后留下的制度与社会后果为其法定安排不等于地方执行，后续须以地域材料追踪负担与成效。桥接条目只标示可回查的年度问题与材料链；实录有中央选择性，崇祯期尤其需要奏疏、地方、朝鲜及满文材料交叉。；本条特别核对报称信息的来源、抄传与行政分类。

\eventanchor{ML-1628-007}{明·1628（崇祯元年）}
\paragraph{1628（崇祯元年）：朝鲜、辽东和海上关系的本年材料提供外部观察位置，明、后金（清）与地方势力的称谓需具体化。（材料核查重点：报称信息的来源、抄传与行政分类）} 此事把权力、法律或文化关系推进到可辨认的节点。账本的证据限度是来源导向的年度桥接条目：本年材料可定位相关政务线索；具体月日、发文机关、地域和执行结果仍须逐项回查，不能以制度文本或奏报替代实际效果。可资核查的材料为《崇祯长编》崇祯元年条及《明史·思宗本纪》；《明史·外国传》；《朝鲜王朝实录》、琉球《历代宝案》或海外档案。它们能够说明制度如何表述自身、某种命令或交涉如何被记录，却未必能直接复原家庭内部关系、性别分工、城市日常或海外行动者的意图。事件发生的条件涉及朝贡名义、边境安全与港口贸易的利益使交涉持续发生；其后留下的制度与社会后果为它为后续册封、互市或海防安排留下文书与跨政权比较线索。桥接条目只标示可回查的年度问题与材料链；实录有中央选择性，崇祯期尤其需要奏疏、地方、朝鲜及满文材料交叉。；本条特别核对报称信息的来源、抄传与行政分类。

\eventanchor{ML-1628-008}{明·1628（崇祯元年）}
\paragraph{1628（崇祯元年）：旱灾、蝗灾、疫病、河患与赈济的本年材料连接环境和社会压力，死亡与流离数字须谨慎处理} 此事把权力、法律或文化关系推进到可辨认的节点。账本的证据限度是来源导向的年度桥接条目：本年材料可定位相关政务线索；具体月日、发文机关、地域和执行结果仍须逐项回查，不能以制度文本或奏报替代实际效果。可资核查的材料为《崇祯长编》崇祯元年条及《明史·思宗本纪》；《明史·五行志》；受灾府县地方志与奏疏。它们能够说明制度如何表述自身、某种命令或交涉如何被记录，却未必能直接复原家庭内部关系、性别分工、城市日常或海外行动者的意图。事件发生的条件涉及自然风险与地方报告促成朝廷或地方的应对记录；其后留下的制度与社会后果为赈济、迁徙和损失的实际范围仍须以受灾地区材料分别核验。桥接条目只标示可回查的年度问题与材料链；实录有中央选择性，崇祯期尤其需要奏疏、地方、朝鲜及满文材料交叉。

\eventanchor{ML-1629-001}{明·1629（崇祯二年）冬}
\paragraph{1629（崇祯二年）冬：己巳事变：后金军队攻破长城，掠夺北京周边地区} 此事把权力、法律或文化关系推进到可辨认的节点。账本的证据限度是编年候选的年序可由官修与后出材料交叉；具体月日、参与者、战果或数字必须回查所列材料，不能将单一叙事当作定论。可资核查的材料为《崇祯长编》崇祯二年条及《明史·思宗本纪》；《明史·兵志》及相关列传；边镇、地方志或对方材料。它们能够说明制度如何表述自身、某种命令或交涉如何被记录，却未必能直接复原家庭内部关系、性别分工、城市日常或海外行动者的意图。事件发生的条件涉及前期边防、地方武装或跨区域资源调动的压力推动了该行动；其后留下的制度与社会后果为它改变了后续调兵、补给或地方秩序的条件；战果和伤亡须分开核对。译写候选以编年材料定位；实录记载反映朝廷选择，崇祯期须另以奏疏、地方及敌对政权材料交叉。

\subsection{家庭、性别与危机材料}

\eventanchor{ML-1629-002}{明·1629（崇祯二年）}
\paragraph{1629（崇祯二年）：崇祯皇帝削减帝国邮政经费，导致山西失业邮政工人叛乱} 此事把权力、法律或文化关系推进到可辨认的节点。账本的证据限度是编年候选的年序可由官修与后出材料交叉；具体月日、参与者、战果或数字必须回查所列材料，不能将单一叙事当作定论。可资核查的材料为《崇祯长编》崇祯二年条及《明史·思宗本纪》；《明会典》相关条目；地方志、黄册/契约/河工材料。它们能够说明制度如何表述自身、某种命令或交涉如何被记录，却未必能直接复原家庭内部关系、性别分工、城市日常或海外行动者的意图。事件发生的条件涉及财政征解、交通工程或货币支付的压力使相关措施进入政务；其后留下的制度与社会后果为其法定安排不等于地方执行，后续须以地域材料追踪负担与成效。译写候选以编年材料定位；实录记载反映朝廷选择，崇祯期须另以奏疏、地方及敌对政权材料交叉。

\eventanchor{ML-1629-003}{明·1629（崇祯二年）}
\paragraph{1629（崇祯二年）：She-An叛乱：叛军被击败} 此事把权力、法律或文化关系推进到可辨认的节点。账本的证据限度是编年候选的年序可由官修与后出材料交叉；具体月日、参与者、战果或数字必须回查所列材料，不能将单一叙事当作定论。可资核查的材料为《崇祯长编》崇祯二年条及《明史·思宗本纪》；《明史·兵志》及相关列传；边镇、地方志或对方材料。它们能够说明制度如何表述自身、某种命令或交涉如何被记录，却未必能直接复原家庭内部关系、性别分工、城市日常或海外行动者的意图。事件发生的条件涉及前期边防、地方武装或跨区域资源调动的压力推动了该行动；其后留下的制度与社会后果为它改变了后续调兵、补给或地方秩序的条件；战果和伤亡须分开核对。译写候选以编年材料定位；实录记载反映朝廷选择，崇祯期须另以奏疏、地方及敌对政权材料交叉。

\eventanchor{ML-1629-004}{明·1629（崇祯二年）}
\paragraph{1629（崇祯二年）：地方结社、宗教、出版或士人文集的本年线索提供战乱中的社会观察，幸存者记忆不能概括全国。（材料核查重点：诏令或奏议的形成与发受链）} 此事把权力、法律或文化关系推进到可辨认的节点。账本的证据限度是来源导向的年度桥接条目：本年材料可定位相关政务线索；具体月日、发文机关、地域和执行结果仍须逐项回查，不能以制度文本或奏报替代实际效果。可资核查的材料为《崇祯长编》崇祯二年条及《明史·思宗本纪》；《明史·选举志》或《艺文志》；文集、碑刻或书目材料。它们能够说明制度如何表述自身、某种命令或交涉如何被记录，却未必能直接复原家庭内部关系、性别分工、城市日常或海外行动者的意图。事件发生的条件涉及制度、教育或知识传播的需要推动了相关行动或文本形成；其后留下的制度与社会后果为它提供制度和传播史的锚点，不能据此推断社会整体接受。桥接条目只标示可回查的年度问题与材料链；实录有中央选择性，崇祯期尤其需要奏疏、地方、朝鲜及满文材料交叉。；本条特别核对诏令或奏议的形成与发受链。

\eventanchor{ML-1629-005}{明·1629（崇祯二年）}
\paragraph{1629（崇祯二年）：辽东防务、后金（清）军事行动和流动作战集团的本年记录标记多战场互动，军数和胜败须保留分歧。（材料核查重点：诏令或奏议的形成与发受链）} 此事把权力、法律或文化关系推进到可辨认的节点。账本的证据限度是来源导向的年度桥接条目：本年材料可定位相关政务线索；具体月日、发文机关、地域和执行结果仍须逐项回查，不能以制度文本或奏报替代实际效果。可资核查的材料为《崇祯长编》崇祯二年条及《明史·思宗本纪》；《明史·兵志》及相关列传；边镇、地方志或对方材料。它们能够说明制度如何表述自身、某种命令或交涉如何被记录，却未必能直接复原家庭内部关系、性别分工、城市日常或海外行动者的意图。事件发生的条件涉及前期边防、地方武装或跨区域资源调动的压力推动了该行动；其后留下的制度与社会后果为它改变了后续调兵、补给或地方秩序的条件；战果和伤亡须分开核对。桥接条目只标示可回查的年度问题与材料链；实录有中央选择性，崇祯期尤其需要奏疏、地方、朝鲜及满文材料交叉。；本条特别核对诏令或奏议的形成与发受链。

\eventanchor{ML-1629-006}{明·1629（崇祯二年）}
\paragraph{1629（崇祯二年）：辽东防务、后金（清）军事行动和流动作战集团的本年记录标记多战场互动，军数和胜败须保留分歧。（材料核查重点：报称信息的来源、抄传与行政分类）} 此事把权力、法律或文化关系推进到可辨认的节点。账本的证据限度是来源导向的年度桥接条目：本年材料可定位相关政务线索；具体月日、发文机关、地域和执行结果仍须逐项回查，不能以制度文本或奏报替代实际效果。可资核查的材料为《崇祯长编》崇祯二年条及《明史·思宗本纪》；《明史·兵志》及相关列传；边镇、地方志或对方材料。它们能够说明制度如何表述自身、某种命令或交涉如何被记录，却未必能直接复原家庭内部关系、性别分工、城市日常或海外行动者的意图。事件发生的条件涉及前期边防、地方武装或跨区域资源调动的压力推动了该行动；其后留下的制度与社会后果为它改变了后续调兵、补给或地方秩序的条件；战果和伤亡须分开核对。桥接条目只标示可回查的年度问题与材料链；实录有中央选择性，崇祯期尤其需要奏疏、地方、朝鲜及满文材料交叉。；本条特别核对报称信息的来源、抄传与行政分类。

\eventanchor{ML-1629-007}{明·1629（崇祯二年）}
\paragraph{1629（崇祯二年）：地方结社、宗教、出版或士人文集的本年线索提供战乱中的社会观察，幸存者记忆不能概括全国。（材料核查重点：报称信息的来源、抄传与行政分类）} 此事把权力、法律或文化关系推进到可辨认的节点。账本的证据限度是来源导向的年度桥接条目：本年材料可定位相关政务线索；具体月日、发文机关、地域和执行结果仍须逐项回查，不能以制度文本或奏报替代实际效果。可资核查的材料为《崇祯长编》崇祯二年条及《明史·思宗本纪》；《明史·选举志》或《艺文志》；文集、碑刻或书目材料。它们能够说明制度如何表述自身、某种命令或交涉如何被记录，却未必能直接复原家庭内部关系、性别分工、城市日常或海外行动者的意图。事件发生的条件涉及制度、教育或知识传播的需要推动了相关行动或文本形成；其后留下的制度与社会后果为它提供制度和传播史的锚点，不能据此推断社会整体接受。桥接条目只标示可回查的年度问题与材料链；实录有中央选择性，崇祯期尤其需要奏疏、地方、朝鲜及满文材料交叉。；本条特别核对报称信息的来源、抄传与行政分类。

\subsection{城市文化与地方知识}

\eventanchor{ML-1629-008}{明·1629（崇祯二年）}
\paragraph{1629（崇祯二年）：天启末至崇祯朝的任免、奏疏和廷议构成本年中枢危机线索；崇祯无实录，必须以多类材料互校} 此事把权力、法律或文化关系推进到可辨认的节点。账本的证据限度是来源导向的年度桥接条目：本年材料可定位相关政务线索；具体月日、发文机关、地域和执行结果仍须逐项回查，不能以制度文本或奏报替代实际效果。可资核查的材料为《崇祯长编》崇祯二年条及《明史·思宗本纪》；《明史·本纪》及相关列传；诏令、奏疏或会典版本。它们能够说明制度如何表述自身、某种命令或交涉如何被记录，却未必能直接复原家庭内部关系、性别分工、城市日常或海外行动者的意图。事件发生的条件涉及继承、官僚协调或皇权运作的压力使问题进入朝廷程序；其后留下的制度与社会后果为它影响后续人事与政策议程；动机和派系标签不能由单一叙事裁定。桥接条目只标示可回查的年度问题与材料链；实录有中央选择性，崇祯期尤其需要奏疏、地方、朝鲜及满文材料交叉。

\eventanchor{ML-1630-001}{明·1630（崇祯三年）夏}
\paragraph{1630（崇祯三年）夏：稷巳之变：后金军队撤退} 此事把权力、法律或文化关系推进到可辨认的节点。账本的证据限度是编年候选的年序可由官修与后出材料交叉；具体月日、参与者、战果或数字必须回查所列材料，不能将单一叙事当作定论。可资核查的材料为《崇祯长编》崇祯三年条及《明史·思宗本纪》；《明史·兵志》及相关列传；边镇、地方志或对方材料。它们能够说明制度如何表述自身、某种命令或交涉如何被记录，却未必能直接复原家庭内部关系、性别分工、城市日常或海外行动者的意图。事件发生的条件涉及前期边防、地方武装或跨区域资源调动的压力推动了该行动；其后留下的制度与社会后果为它改变了后续调兵、补给或地方秩序的条件；战果和伤亡须分开核对。译写候选以编年材料定位；实录记载反映朝廷选择，崇祯期须另以奏疏、地方及敌对政权材料交叉。

\eventanchor{ML-1630-002}{明·1630（崇祯三年）}
\paragraph{1630（崇祯三年）：朝鲜、辽东和海上关系的本年材料提供外部观察位置，明、后金（清）与地方势力的称谓需具体化。（材料核查重点：诏令或奏议的形成与发受链）} 此事把权力、法律或文化关系推进到可辨认的节点。账本的证据限度是来源导向的年度桥接条目：本年材料可定位相关政务线索；具体月日、发文机关、地域和执行结果仍须逐项回查，不能以制度文本或奏报替代实际效果。可资核查的材料为《崇祯长编》崇祯三年条及《明史·思宗本纪》；《明史·外国传》；《朝鲜王朝实录》、琉球《历代宝案》或海外档案。它们能够说明制度如何表述自身、某种命令或交涉如何被记录，却未必能直接复原家庭内部关系、性别分工、城市日常或海外行动者的意图。事件发生的条件涉及朝贡名义、边境安全与港口贸易的利益使交涉持续发生；其后留下的制度与社会后果为它为后续册封、互市或海防安排留下文书与跨政权比较线索。桥接条目只标示可回查的年度问题与材料链；实录有中央选择性，崇祯期尤其需要奏疏、地方、朝鲜及满文材料交叉。；本条特别核对诏令或奏议的形成与发受链。

\eventanchor{ML-1630-003}{明·1630（崇祯三年）}
\paragraph{1630（崇祯三年）：旱灾、蝗灾、疫病、河患与赈济的本年材料连接环境和社会压力，死亡与流离数字须谨慎处理。（材料核查重点：诏令或奏议的形成与发受链）} 此事把权力、法律或文化关系推进到可辨认的节点。账本的证据限度是来源导向的年度桥接条目：本年材料可定位相关政务线索；具体月日、发文机关、地域和执行结果仍须逐项回查，不能以制度文本或奏报替代实际效果。可资核查的材料为《崇祯长编》崇祯三年条及《明史·思宗本纪》；《明史·五行志》；受灾府县地方志与奏疏。它们能够说明制度如何表述自身、某种命令或交涉如何被记录，却未必能直接复原家庭内部关系、性别分工、城市日常或海外行动者的意图。事件发生的条件涉及自然风险与地方报告促成朝廷或地方的应对记录；其后留下的制度与社会后果为赈济、迁徙和损失的实际范围仍须以受灾地区材料分别核验。桥接条目只标示可回查的年度问题与材料链；实录有中央选择性，崇祯期尤其需要奏疏、地方、朝鲜及满文材料交叉。；本条特别核对诏令或奏议的形成与发受链。

\eventanchor{ML-1630-004}{明·1630（崇祯三年）}
\paragraph{1630（崇祯三年）：旱灾、蝗灾、疫病、河患与赈济的本年材料连接环境和社会压力，死亡与流离数字须谨慎处理。（材料核查重点：报称信息的来源、抄传与行政分类）} 此事把权力、法律或文化关系推进到可辨认的节点。账本的证据限度是来源导向的年度桥接条目：本年材料可定位相关政务线索；具体月日、发文机关、地域和执行结果仍须逐项回查，不能以制度文本或奏报替代实际效果。可资核查的材料为《崇祯长编》崇祯三年条及《明史·思宗本纪》；《明史·五行志》；受灾府县地方志与奏疏。它们能够说明制度如何表述自身、某种命令或交涉如何被记录，却未必能直接复原家庭内部关系、性别分工、城市日常或海外行动者的意图。事件发生的条件涉及自然风险与地方报告促成朝廷或地方的应对记录；其后留下的制度与社会后果为赈济、迁徙和损失的实际范围仍须以受灾地区材料分别核验。桥接条目只标示可回查的年度问题与材料链；实录有中央选择性，崇祯期尤其需要奏疏、地方、朝鲜及满文材料交叉。；本条特别核对报称信息的来源、抄传与行政分类。

\eventanchor{ML-1630-005}{明·1630（崇祯三年）}
\paragraph{1630（崇祯三年）：朝鲜、辽东和海上关系的本年材料提供外部观察位置，明、后金（清）与地方势力的称谓需具体化。（材料核查重点：报称信息的来源、抄传与行政分类）} 此事把权力、法律或文化关系推进到可辨认的节点。账本的证据限度是来源导向的年度桥接条目：本年材料可定位相关政务线索；具体月日、发文机关、地域和执行结果仍须逐项回查，不能以制度文本或奏报替代实际效果。可资核查的材料为《崇祯长编》崇祯三年条及《明史·思宗本纪》；《明史·外国传》；《朝鲜王朝实录》、琉球《历代宝案》或海外档案。它们能够说明制度如何表述自身、某种命令或交涉如何被记录，却未必能直接复原家庭内部关系、性别分工、城市日常或海外行动者的意图。事件发生的条件涉及朝贡名义、边境安全与港口贸易的利益使交涉持续发生；其后留下的制度与社会后果为它为后续册封、互市或海防安排留下文书与跨政权比较线索。桥接条目只标示可回查的年度问题与材料链；实录有中央选择性，崇祯期尤其需要奏疏、地方、朝鲜及满文材料交叉。；本条特别核对报称信息的来源、抄传与行政分类。

\subsection{海外关系与政权观察}

\eventanchor{ML-1630-006}{明·1630（崇祯三年）}
\paragraph{1630（崇祯三年）：朝鲜、辽东和海上关系的本年材料提供外部观察位置，明、后金（清）与地方势力的称谓需具体化。（材料核查重点：同年地方材料所见的执行范围）} 此事把权力、法律或文化关系推进到可辨认的节点。账本的证据限度是来源导向的年度桥接条目：本年材料可定位相关政务线索；具体月日、发文机关、地域和执行结果仍须逐项回查，不能以制度文本或奏报替代实际效果。可资核查的材料为《崇祯长编》崇祯三年条及《明史·思宗本纪》；《明史·外国传》；《朝鲜王朝实录》、琉球《历代宝案》或海外档案。它们能够说明制度如何表述自身、某种命令或交涉如何被记录，却未必能直接复原家庭内部关系、性别分工、城市日常或海外行动者的意图。事件发生的条件涉及朝贡名义、边境安全与港口贸易的利益使交涉持续发生；其后留下的制度与社会后果为它为后续册封、互市或海防安排留下文书与跨政权比较线索。桥接条目只标示可回查的年度问题与材料链；实录有中央选择性，崇祯期尤其需要奏疏、地方、朝鲜及满文材料交叉。；本条特别核对同年地方材料所见的执行范围。

\eventanchor{ML-1630-007}{明·1630（崇祯三年）}
\paragraph{1630（崇祯三年）：旱灾、蝗灾、疫病、河患与赈济的本年材料连接环境和社会压力，死亡与流离数字须谨慎处理。（材料核查重点：同年地方材料所见的执行范围）} 此事把权力、法律或文化关系推进到可辨认的节点。账本的证据限度是来源导向的年度桥接条目：本年材料可定位相关政务线索；具体月日、发文机关、地域和执行结果仍须逐项回查，不能以制度文本或奏报替代实际效果。可资核查的材料为《崇祯长编》崇祯三年条及《明史·思宗本纪》；《明史·五行志》；受灾府县地方志与奏疏。它们能够说明制度如何表述自身、某种命令或交涉如何被记录，却未必能直接复原家庭内部关系、性别分工、城市日常或海外行动者的意图。事件发生的条件涉及自然风险与地方报告促成朝廷或地方的应对记录；其后留下的制度与社会后果为赈济、迁徙和损失的实际范围仍须以受灾地区材料分别核验。桥接条目只标示可回查的年度问题与材料链；实录有中央选择性，崇祯期尤其需要奏疏、地方、朝鲜及满文材料交叉。；本条特别核对同年地方材料所见的执行范围。

\eventanchor{ML-1631-001}{明·1631（崇祯四年）四月}
\paragraph{1631（崇祯四年）四月：叛军占领平度} 此事把权力、法律或文化关系推进到可辨认的节点。账本的证据限度是编年候选的年序可由官修与后出材料交叉；具体月日、参与者、战果或数字必须回查所列材料，不能将单一叙事当作定论。可资核查的材料为《崇祯长编》崇祯四年条及《明史·思宗本纪》；《明史·兵志》及相关列传；边镇、地方志或对方材料。它们能够说明制度如何表述自身、某种命令或交涉如何被记录，却未必能直接复原家庭内部关系、性别分工、城市日常或海外行动者的意图。事件发生的条件涉及前期边防、地方武装或跨区域资源调动的压力推动了该行动；其后留下的制度与社会后果为它改变了后续调兵、补给或地方秩序的条件；战果和伤亡须分开核对。译写候选以编年材料定位；实录记载反映朝廷选择，崇祯期须另以奏疏、地方及敌对政权材料交叉。

\eventanchor{ML-1631-002}{明·1631（崇祯四年）十一月21日}
\paragraph{1631（崇祯四年）十一月21日：大凌河之战：后金夺取大凌河} 此事把权力、法律或文化关系推进到可辨认的节点。账本的证据限度是编年候选的年序可由官修与后出材料交叉；具体月日、参与者、战果或数字必须回查所列材料，不能将单一叙事当作定论。可资核查的材料为《崇祯长编》崇祯四年条及《明史·思宗本纪》；《明史·兵志》及相关列传；边镇、地方志或对方材料。它们能够说明制度如何表述自身、某种命令或交涉如何被记录，却未必能直接复原家庭内部关系、性别分工、城市日常或海外行动者的意图。事件发生的条件涉及前期边防、地方武装或跨区域资源调动的压力推动了该行动；其后留下的制度与社会后果为它改变了后续调兵、补给或地方秩序的条件；战果和伤亡须分开核对。译写候选以编年材料定位；实录记载反映朝廷选择，崇祯期须另以奏疏、地方及敌对政权材料交叉。

\eventanchor{ML-1631-003}{明·1631（崇祯四年）}
\paragraph{1631（崇祯四年）：天启末至崇祯朝的任免、奏疏和廷议构成本年中枢危机线索；崇祯无实录，必须以多类材料互校。（材料核查重点：诏令或奏议的形成与发受链）} 此事把权力、法律或文化关系推进到可辨认的节点。账本的证据限度是来源导向的年度桥接条目：本年材料可定位相关政务线索；具体月日、发文机关、地域和执行结果仍须逐项回查，不能以制度文本或奏报替代实际效果。可资核查的材料为《崇祯长编》崇祯四年条及《明史·思宗本纪》；《明史·本纪》及相关列传；诏令、奏疏或会典版本。它们能够说明制度如何表述自身、某种命令或交涉如何被记录，却未必能直接复原家庭内部关系、性别分工、城市日常或海外行动者的意图。事件发生的条件涉及继承、官僚协调或皇权运作的压力使问题进入朝廷程序；其后留下的制度与社会后果为它影响后续人事与政策议程；动机和派系标签不能由单一叙事裁定。桥接条目只标示可回查的年度问题与材料链；实录有中央选择性，崇祯期尤其需要奏疏、地方、朝鲜及满文材料交叉。；本条特别核对诏令或奏议的形成与发受链。

\eventanchor{ML-1631-004}{明·1631（崇祯四年）}
\paragraph{1631（崇祯四年）：天启末至崇祯朝的任免、奏疏和廷议构成本年中枢危机线索；崇祯无实录，必须以多类材料互校。（材料核查重点：报称信息的来源、抄传与行政分类）} 此事把权力、法律或文化关系推进到可辨认的节点。账本的证据限度是来源导向的年度桥接条目：本年材料可定位相关政务线索；具体月日、发文机关、地域和执行结果仍须逐项回查，不能以制度文本或奏报替代实际效果。可资核查的材料为《崇祯长编》崇祯四年条及《明史·思宗本纪》；《明史·本纪》及相关列传；诏令、奏疏或会典版本。它们能够说明制度如何表述自身、某种命令或交涉如何被记录，却未必能直接复原家庭内部关系、性别分工、城市日常或海外行动者的意图。事件发生的条件涉及继承、官僚协调或皇权运作的压力使问题进入朝廷程序；其后留下的制度与社会后果为它影响后续人事与政策议程；动机和派系标签不能由单一叙事裁定。桥接条目只标示可回查的年度问题与材料链；实录有中央选择性，崇祯期尤其需要奏疏、地方、朝鲜及满文材料交叉。；本条特别核对报称信息的来源、抄传与行政分类。
